% !PRINTED 647
\weberpage{647}
Daraus folgt:

2. \gesperrt{Ist ein Produkt $\alpha\beta$ durch $\pi$ teilbar, so ist wenigstens einer der beiden Faktoren, $\alpha$ oder $\beta$, durch $\pi$ teilbar.}

Denn ist weder $\alpha$ noch $\beta$ durch $\pi$ teilbar, so sind beide relativ prim zu $\pi$ und nach \S\,181, 3. ist $\pi$ auch relativ prim zu $\alpha\beta$.

3. \gesperrt{Jedes von Null verschiedene ganze Funktional $\overline{\omega}$ ist durch ein Primfunktional teilbar.}

Denn ist $\overline{\omega}$ nicht selbst ein Primfunktional, so ist es durch ein von $\omega$ verschiedenes Funktional $\alpha$ teilbar. % [sic] printed without the expected bar over omega
Ist
\[
\overline{\omega} = \alpha \overline{\omega}',
\]
so ist $\overline{\omega}'$ keine Einheit und
\[
N_a(\overline{\omega}) = N_a(\alpha)\, N_a(\overline{\omega}').
\]

Die hier vorkommenden absoluten Normen sind ganze Funktionen in $z$ und der Grad von $N_a(\overline{\omega}')$ ist kleiner als der Grad von $N_a(\omega)$. % [sic] printed without the expected bar over omega
Ist $\overline{\omega}'$ noch kein Primfunktional, so kann man so fortfahren und muß schließlich auf einen Primfaktor $\pi$ von $\overline{\omega}$ kommen.

In dieser Weise schließt man weiter wie in Bd.\,II, \S\,158, wobei nur an Stelle der dort benutzten \gesperrt{Größe} der ganzen rationalen Zahlen hier die \gesperrt{Höhe des Grades} ganzer rationaler Funktionen der Variablen $z$ tritt. So erhält man auch die Sätze:

4. \gesperrt{Jedes ganze Funktional $\overline{\omega}$ in $\overline{\Omega}$, das keine Einheit ist, läßt sich in eine endliche Anzahl von Primfaktoren zerlegen, und zwar nur auf eine Weise, wenn assoziierte Funktionale als nicht verschieden betrachtet werden.}

Im folgenden müssen wir von dem Satz Gebrauch machen, daß eine ganze rationale Funktion von $z$ mit numerischen Koeffizienten in lineare Faktoren zerlegbar ist, mit anderen Worten, wir müssen den Fundamentalsatz der Algebra von der Wurzelexistenz voraussetzen. Es folgt daraus zunächst:

5. \gesperrt{Die ganze rationale Funktion von $z$ niedrigsten Grades, die durch ein Primfunktional $\pi$ teilbar ist, ist eine lineare Funktion $z - c$.}

Nach \S\,180, 9. gibt es überhaupt ganz % [sic] printed without final -e
rationale Funktionen von $z$, die durch $\pi$ teilbar sind. Zerlegt man eine solche Funk-
%CONT>>
